\documentclass[10pt]{article}

\usepackage[margin=0.5in]{geometry}
\usepackage[T1]{fontenc}
\usepackage{lmodern}
\usepackage{amsmath}
\usepackage{xcolor}
\usepackage{listings}
\usepackage{booktabs}
\usepackage{array}
\usepackage{multicol}
\usepackage{enumitem}
\usepackage[hidelinks]{hyperref}

\setlength{\parindent}{0pt}
\setlength{\parskip}{2pt}
\setlength{\columnsep}{16pt}
\setlist{nosep, leftmargin=1.1em}

\definecolor{kw}{RGB}{0,70,160}
\definecolor{cm}{RGB}{95,125,95}
\definecolor{st}{RGB}{165,60,30}
\definecolor{codebg}{RGB}{247,247,249}

\lstdefinestyle{py}{
  language=Python,
  basicstyle=\ttfamily\footnotesize,
  keywordstyle=\color{kw}\bfseries,
  commentstyle=\color{cm}\itshape,
  stringstyle=\color{st},
  morekeywords={nonlocal,True,False,None,yield},
  showstringspaces=false,
  columns=fullflexible,
  keepspaces=true,
  breaklines=true,
  upquote=true,
  tabsize=4,
  frame=single,
  rulecolor=\color{black!12},
  backgroundcolor=\color{codebg},
  xleftmargin=2pt, xrightmargin=2pt,
  aboveskip=3pt, belowskip=3pt
}
\lstnewenvironment{py}{\lstset{style=py}}{}
\lstnewenvironment{tpl}{\lstset{style=py}}{}   % self-contained templates (extracted + tested)

% inline code: never put the characters # % \ or unbalanced braces inside \cd{}
\newcommand{\cd}[1]{\texttt{\detokenize{#1}}}

\makeatletter
\renewcommand\section{\@startsection{section}{1}{0pt}{6pt plus 2pt}{3pt}{\large\bfseries\color{kw}}}
\renewcommand\subsection{\@startsection{subsection}{2}{0pt}{4pt plus 1pt}{2pt}{\normalsize\bfseries}}
\makeatother

\begin{document}

{\centering {\LARGE\bfseries Python Review Cheatsheet}\\[2pt]
{\small Fundamentals, standard library, performance, and algorithm patterns for timed coding assessments}\par}
\vspace{4pt}

\begin{multicols}{2}
\small

% =====================================================================
\section{Complexity budget}
CPython does about $10^7$ simple operations per second; budget about $10^7$--$5\cdot10^7$ per test.

\begin{tabular}{@{}ll@{}}
\toprule
Max $n$ & Target complexity \\
\midrule
$\le 10$ & $O(n!)$ \\
$\le 20$ & $O(2^n)$, $O(n 2^n)$ \\
$\le 500$ & $O(n^3)$ (tight) \\
$\le 5\cdot 10^3$ & $O(n^2)$ \\
$\le 10^6$ & $O(n \log n)$, $O(n)$ \\
$\ge 10^9$ & $O(\log n)$, $O(1)$, math \\
\bottomrule
\end{tabular}

Grids: touch each cell $O(1)$ (or $O(\log)$) times. Values up to $10^{12}$ mean you can't simulate step by step; use arithmetic.

% =====================================================================
\section{Numbers}
\begin{py}
7 // 2, -7 // 2        # 3, -4   floor toward -inf
7 % 3, -7 % 3          # 1, 2    sign of divisor
int(-7 / 2)            # -3      truncate toward 0
divmod(17, 5)          # (3, 2)
-(-a // b)             # ceil(a / b) for ints, exact
2 ** 10, pow(2, 10, MOD)    # fast modular power
pow(a, -1, MOD)        # modular inverse (3.8+)
abs(x), round(2.5)     # 2 !  banker's rounding
int(x + 0.5)           # round half up (x >= 0)
float('inf'), -math.inf
10**18 + 1             # ints never overflow
0.1 + 0.2 == 0.3       # False -> math.isclose
int("ff", 16), int("101", 2), bin(5), hex(255)
1_000_000              # digit separators
\end{py}

\subsection{Bit tricks}
\begin{py}
x & 1                  # odd?
x >> k & 1             # k-th bit
x | (1 << k)           # set bit k
x & ~(1 << k)          # clear bit k
x ^ (1 << k)           # toggle bit k
x & (x - 1)            # drop lowest set bit
x & -x                 # lowest set bit
x.bit_count()          # popcount (3.10+); bin(x).count("1")
x.bit_length()         # floor(log2 x) + 1
x & (x - 1) == 0       # power of two (x > 0)
a ^ b ^ a == b         # XOR cancels pairs
for sub in range(1 << n): ...           # all subsets
s = mask
while s:                                 # all submasks
    s = (s - 1) & mask
\end{py}

% =====================================================================
\section{Core syntax}
\begin{py}
a, b = b, a                       # swap
first, *rest = lst                # unpacking
x = a if cond else b              # ternary
if (n := len(s)) > 10: ...        # walrus
0 <= r < R and 0 <= c < C         # chained compare
for i, x in enumerate(a, start=1): ...
for x, y in zip(a, b): ...        # shortest wins
for x, y in zip(a, a[1:]): ...    # adjacent pairs
for i in reversed(range(n)): ...
for i in range(n - 1, -1, -1): ...
while True: ... break
for ...: ... else: ...            # else runs if no break
\end{py}

\subsection{Truthiness and identity}
Falsy: \cd{0, 0.0, "", [], {}, set(), None, False}. Use \cd{is None} / \cd{is not None} when 0 is a valid
value. \cd{==} compares values; \cd{is} compares identity.

\subsection{Comprehensions and generators}
\begin{py}
[f(x) for x in it if p(x)]
[x if p(x) else y for x in it]   # map with ternary
{k: v for k, v in pairs}
{x % 10 for x in nums}
sum(x * x for x in nums)          # generator: no list
[[0] * C for _ in range(R)]       # 2D grid (correct)
[v for row in grid for v in row]  # flatten
def gen():                        # lazy sequence
    yield 1
    yield from range(3)
\end{py}

\subsection{Functions and scope}
\begin{py}
def f(a, b=0, *args, key=None, **kw): ...
lambda x: x[1]
def outer():
    cnt = 0
    def inner():
        nonlocal cnt              # needed to REBIND
        cnt += 1                  # (mutating a list is fine without it)
\end{py}
Names assigned inside a function are local unless declared \cd{nonlocal}/\cd{global}.

\subsection{Classes (for nodes, heap items)}
\begin{py}
class Node:
    __slots__ = ("val", "next")   # faster, less memory
    def __init__(self, val, nxt=None):
        self.val, self.next = val, nxt
    def __lt__(self, other):      # makes it heap-comparable
        return self.val < other.val
    def __repr__(self):
        return f"Node({self.val})"

from dataclasses import dataclass, field
@dataclass(order=True)
class Item:
    prio: int
    name: str = field(compare=False)
\end{py}

% =====================================================================
\section{Strings}
\begin{py}
s[i], s[-1], s[a:b], s[::-1]
"".join(parts)          # O(total); s += c in a loop may be O(n^2)
s.split(), s.split(","), s.splitlines()
s.strip(), s.lstrip("0"), s.rstrip()
s.lower(), s.upper(), s.swapcase(), s.title()
s.startswith(p), s.endswith(p)
s.find(t)               # -1 if missing; s.index(t) raises
s.count(t), s.replace(a, b), t in s
s.isdigit(), s.isalpha(), s.isalnum(), s.isspace()
ord("a"), chr(97); ord(c) - ord("a")     # 0..25
s.zfill(4), s.rjust(5), s.ljust(5, "."), s.center(9)
f"{x:.2f} {n:05d} {s:>8} {v!r} {big:,}"
"-".join(map(str, nums))
list(s) / "".join(chars)                 # strings are immutable
s[:i] + c + s[i + 1:]                    # replace one char (O(n))
import string; string.ascii_lowercase, string.digits
sorted(s), "".join(sorted(s))            # anagram key
\end{py}

% =====================================================================
\section{Lists and tuples}
\begin{py}
a.append(x); a.extend(it); a.insert(i, x)
a.pop(); a.pop(i); a.remove(x); del a[i]; del a[i:j]
a.index(x); a.count(x); a.clear()
a.sort(); a.reverse()         # in place, return None
a[::-1], a[::2], a[1:-1]      # slices are copies
a[i:j] = [x, y]               # slice assignment
a * 3, a + b                  # new lists
b = a[:]; b = list(a); b = a.copy()       # shallow copy
g2 = [row[:] for row in grid]             # copy a grid
import copy; copy.deepcopy(x)             # slow
max(a, key=f), min(a, default=0)
sum(a), any(a), all(a), len(a)
a.sort(key=lambda t: (-t[1], t[0]))       # multi-key
t = (1, 2); t + (3,); (x,)                 # 1-tuple needs comma
\end{py}
\textbf{Matrix tricks:}
\begin{py}
list(zip(*m))                     # transpose (tuples)
[list(r) for r in zip(*m[::-1])]  # rotate 90 cw
[list(r) for r in zip(*m)][::-1]  # rotate 90 ccw
[row[::-1] for row in m]          # mirror left-right
m[::-1]                           # flip upside-down
# diagonals: r - c is constant on "\" diagonals, r + c on "/"
# flatten index: k = r * C + c;  r, c = divmod(k, C)
\end{py}

% =====================================================================
\section{dict / set / collections}
\begin{py}
d.get(k, 0); d.setdefault(k, []).append(v)
d.pop(k, None); k in d; d.items(); d.keys()
d | other; d.update(other)          # merge
dict(zip(keys, vals)); dict.fromkeys(keys, 0)
max(d, key=d.get)                   # key with max value
# insertion order is preserved (3.7+)

s.add(x); s.discard(x); s.remove(x); s.pop()
a | b, a & b, a - b, a ^ b, a <= b, a.isdisjoint(b)
frozenset(cells)        # hashable set
tuple(lst)              # hashable list
# hashable: int, str, tuple(of hashables), frozenset

from collections import defaultdict, Counter, deque, OrderedDict
g = defaultdict(list); g[u].append(v)
# reading g[x] INSERTS x -> use `x in g` to test
c = Counter(it); c[x] += 1; c[missing] == 0
c.most_common(k); c.total()         # total: 3.10+
c1 + c2, c1 - c2, c1 & c2, c1 | c2  # multiset ops
Counter(a) == Counter(b)            # anagram check
od = OrderedDict(); od.move_to_end(k)
od.popitem(last=False)              # pop oldest -> LRU cache
\end{py}

\subsection{deque}
\begin{py}
q = deque([x]); q.append(x); q.appendleft(x)
q.pop(); q.popleft(); q[0]; q[-1]   # O(1) ends
q.rotate(k); deque(maxlen=k)        # ring buffer
\end{py}

\subsection{heapq (min-heap on a list)}
\begin{py}
import heapq
heapq.heappush(h, (prio, i, item))  # i breaks ties
prio, _, item = heapq.heappop(h)
h[0]                                # peek
heapq.heapify(a)                    # O(n)
heapq.heappushpop(h, x); heapq.heapreplace(h, x)
heapq.nlargest(k, it, key=f); heapq.nsmallest(k, it)
heapq.merge(*sorted_lists)          # lazy k-way merge
heapq.heappush(h, -x)               # max-heap via negation
\end{py}

\subsection{bisect (sorted lists)}
\begin{py}
import bisect
bisect.bisect_left(a, x)    # first i with a[i] >= x
bisect.bisect_right(a, x)   # first i with a[i] >  x
bisect.insort(a, x)         # O(n) insert
bisect.bisect_left(a, x, key=f)     # key= is 3.10+
n_in = bisect.bisect_right(a, hi) - bisect.bisect_left(a, lo)
\end{py}

% =====================================================================
\section{Sorting}
\begin{py}
sorted(a, key=f, reverse=True)       # stable, O(n log n)
a.sort(key=lambda x: (len(x), x))
from operator import itemgetter, attrgetter
a.sort(key=itemgetter(1, 0))
order = sorted(range(n), key=a.__getitem__)   # argsort
pos = [0] * n
for rank, i in enumerate(order):
    pos[i] = rank                    # inverse permutation
from functools import cmp_to_key
a.sort(key=cmp_to_key(lambda x, y: (x > y) - (x < y)))
\end{py}
Stability lets you sort by secondary key first, then primary.

% =====================================================================
\section{itertools / functools / math}
\begin{py}
from itertools import (accumulate, product, permutations,
    combinations, combinations_with_replacement, groupby,
    pairwise, chain, count, cycle, islice, repeat, starmap)
list(accumulate(a, initial=0))       # prefix sums
list(accumulate(a, max))             # running max
product(range(3), repeat=2)          # nested loops
permutations(a, r); combinations(a, r)
[(k, len(list(g))) for k, g in groupby(s)]   # RLE
chain.from_iterable(lists)           # flatten one level
islice(it, 10)

from functools import cache, lru_cache, reduce, partial
@cache                               # memoize (args hashable)
def f(i, j): ...
f.cache_clear()
reduce(lambda x, y: x * y, a, 1); math.prod(a)

import math
math.gcd(a, b), math.lcm(a, b), math.isqrt(n)
math.comb(n, k), math.perm(n, k), math.factorial(n)
math.sqrt, math.log(x, b), math.log2, math.hypot
math.floor, math.ceil, math.inf, math.pi
math.isclose(a, b, rel_tol=1e-9, abs_tol=1e-9)
\end{py}

% =====================================================================
\section{Recursion}
\begin{py}
import sys
sys.setrecursionlimit(10**6)   # default ~1000 frames
\end{py}
\begin{itemize}
  \item Fine for depth $\le$ a few thousand: backtracking, balanced trees, memoised DP on small ranges.
  \item Grids, chains and graphs of size $\ge 10^4$: use an explicit stack/queue. Very deep recursion can
        crash the interpreter even with a raised limit.
  \item Memoised recursion = top-down DP; convert to bottom-up if depth or speed is a problem.
\end{itemize}

% =====================================================================
\section{Operation costs}
\begin{tabular}{@{}p{1.4cm}p{3.1cm}p{3.3cm}@{}}
\toprule
Type & $O(1)$ & $O(n)$ / other \\
\midrule
list & \cd{a[i]}, append, pop(), len & \cd{x in a}, insert/pop(0), remove, slice, \cd{a+b}, index \\
dict/set & get/set/del, \cd{in}, add & copy, iterate \\
deque & both ends & middle index \\
heap & \cd{h[0]} & push/pop $\log n$; heapify $n$ \\
str & index, len & slice, \cd{+}, \cd{in}, join \\
sort & -- & $n \log n$ \\
\bottomrule
\end{tabular}

% =====================================================================
\section{Performance tricks}
\begin{itemize}
  \item Convert to a \cd{set} before repeated membership tests.
  \item Dense integer keys: use a list, not a dict. Flatten \cd{(r, c)} to \cd{r * C + c}.
  \item Built-ins (\cd{sum}, \cd{max}, \cd{sorted}, slicing, \cd{str.join}, \cd{Counter}) run in C. Prefer them over manual loops.
  \item Hoist lookups: \cd{row = grid[r]}, \cd{append = res.append}, local variables over globals.
  \item Carry running totals instead of recomputing \cd{sum(path)}.
  \item Avoid \cd{deepcopy}, string \cd{+=} in loops, and \cd{list.pop(0)}.
  \item Big integers as bitsets: \cd{bits |= bits << x} gives subset sums fast.
  \item Fast input: \cd{data = sys.stdin.buffer.read().split()}.
  \item Avoid \cd{print} in hot loops.
\end{itemize}

% =====================================================================
\section{Pitfalls}
\begin{itemize}
  \item \cd{[[0] * C] * R} aliases rows. Use a comprehension.
  \item Mutable default args (\cd{def f(a=[])}) are shared across calls.
  \item \cd{res.append(path)} stores a reference. Append \cd{path[:]}.
  \item Lambdas in loops capture the variable, not its value: \cd{lambda i=i: i}.
  \item \cd{a.sort()} returns \cd{None}; don't write \cd{a = a.sort()}.
  \item Modifying a list/dict while iterating over it. Iterate over a copy or build a new one.
  \item \cd{defaultdict} lookup inserts keys (it changes \cd{len} and iteration).
  \item \cd{round()} is banker's rounding; \cd{/} always returns a float.
  \item \cd{min()}/\cd{max()} on an empty sequence raises. Pass \cd{default=}.
  \item \cd{str.split()} vs \cd{str.split(" ")} handle repeated spaces differently.
  \item \cd{@cache} on a nested function is rebuilt on every outer call (good); on a global one it persists across tests (call \cd{cache_clear}).
  \item Shadowing built-ins: \cd{sum}, \cd{list}, \cd{max}, \cd{id}, \cd{input}, \cd{map}.
  \item Integer vs.\ float: never compare floats with \cd{==}; prefer integer arithmetic.
\end{itemize}

% =====================================================================
\section{Pattern recognition}
\begin{tabular}{@{}p{3.5cm}p{4.3cm}@{}}
\toprule
Signal in the problem & Reach for \\
\midrule
fewest moves / steps, unweighted & BFS (level by level) \\
many starting points & multi-source BFS \\
states are strings/tuples & implicit-graph BFS \\
weighted shortest path & Dijkstra (heap) \\
connected groups, regions & DFS/BFS or union--find \\
prerequisites, ordering & topological sort \\
each node has one next & functional graph walk \\
all combinations / arrangements & backtracking \\
count ways / best value, overlapping subproblems & DP \\
two players, optimal play & minimax interval DP \\
contiguous subarray sum & prefix sums (+ hashmap) \\
range sums in a grid & 2D prefix sums \\
longest/shortest window with a property & sliding window \\
sorted input, pair/triple sums & two pointers \\
max/min $x$ such that ok($x$), monotone & binary search on answer \\
next greater/smaller element & monotonic stack \\
top $k$, running median, merge $k$ & heap(s) \\
overlapping intervals & sort by start (merge) / end (select) \\
reachability, local choice suffices & greedy (prove or brute-check) \\
huge step counts of shifts & track cumulative offset / extremes \\
synchronous grid update & new grid or bit-encoded next state \\
$n \le 20$ & bitmask / brute force \\
\bottomrule
\end{tabular}

% =====================================================================
\section{Templates: graphs and grids}

\subsection{Grid helpers}
\begin{py}
R, C = len(g), len(g[0])
D4 = [(1, 0), (-1, 0), (0, 1), (0, -1)]
D8 = [(a, b) for a in (-1, 0, 1) for b in (-1, 0, 1) if a or b]
CW = [(-1, 0), (0, 1), (1, 0), (0, -1)]   # up right down left
d = (d + 1) % 4   # turn right;  (d - 1) % 4 turn left
for dr, dc in D4:
    nr, nc = r + dr, c + dc
    if 0 <= nr < R and 0 <= nc < C: ...
\end{py}

\subsection{BFS (shortest path, unweighted)}
\begin{tpl}
from collections import deque

def bfs(start, target, neighbors):
    q, seen, steps = deque([start]), {start}, 0
    while q:
        for _ in range(len(q)):          # one level
            u = q.popleft()
            if u == target:
                return steps
            for v in neighbors(u):
                if v not in seen:        # mark on ENQUEUE
                    seen.add(v)
                    q.append(v)
        steps += 1
    return -1
\end{tpl}
Multi-source: start the queue with every source. 0--1 BFS (edge weights 0/1): deque, \cd{appendleft}
for 0-weight edges. Bidirectional BFS: expand the smaller frontier.

\subsection{Iterative DFS / flood fill}
\begin{tpl}
def flood(grid, r, c, seen):
    R, C = len(grid), len(grid[0])
    stack, cells = [(r, c)], []
    seen[r][c] = True
    while stack:
        a, b = stack.pop()
        cells.append((a, b))
        for x, y in ((a+1, b), (a-1, b), (a, b+1), (a, b-1)):
            if (0 <= x < R and 0 <= y < C and not seen[x][y]
                    and grid[x][y] == grid[r][c]):
                seen[x][y] = True
                stack.append((x, y))
    return cells
\end{tpl}

\subsection{Dijkstra}
\begin{tpl}
import heapq

def dijkstra(adj, src):              # adj[u] = [(v, w), ...]
    dist = {src: 0}
    h = [(0, src)]
    while h:
        d, u = heapq.heappop(h)
        if d > dist[u]:
            continue                 # stale entry
        for v, w in adj[u]:
            nd = d + w
            if nd < dist.get(v, float("inf")):
                dist[v] = nd
                heapq.heappush(h, (nd, v))
    return dist
\end{tpl}

\subsection{Union--Find}
\begin{tpl}
class DSU:
    def __init__(self, n):
        self.p, self.sz = list(range(n)), [1] * n
    def find(self, x):
        while self.p[x] != x:
            self.p[x] = self.p[self.p[x]]   # path halving
            x = self.p[x]
        return x
    def union(self, a, b):
        a, b = self.find(a), self.find(b)
        if a == b:
            return False                    # already joined
        if self.sz[a] < self.sz[b]:
            a, b = b, a
        self.p[b] = a
        self.sz[a] += self.sz[b]
        return True
\end{tpl}

\subsection{Topological sort / cycle detection}
\begin{tpl}
from collections import deque

def topo_sort(n, edges):             # (u, v): u before v
    adj, indeg = [[] for _ in range(n)], [0] * n
    for u, v in edges:
        adj[u].append(v)
        indeg[v] += 1
    q = deque(i for i in range(n) if indeg[i] == 0)
    order = []
    while q:
        u = q.popleft()
        order.append(u)
        for v in adj[u]:
            indeg[v] -= 1
            if indeg[v] == 0:
                q.append(v)
    return order if len(order) == n else None   # None: cycle
\end{tpl}
DFS alternative: colours 0 = new, 1 = on path, 2 = done; an edge to colour 1 is a cycle. Undirected
cycle: union--find returns \cd{False}, or DFS that ignores the parent edge.

\subsection{Functional graph (one outgoing edge each)}
\begin{py}
seen, cur = {start}, start
while cur in nxt and nxt[cur] not in seen:
    cur = nxt[cur]; seen.add(cur)
# Floyd (O(1) memory): slow = f(slow); fast = f(f(fast))
\end{py}

\subsection{Trees}
\begin{py}
def depth(node):                      # postorder: combine children
    if not node: return 0
    return 1 + max(depth(node.left), depth(node.right))
# level order: deque, `for _ in range(len(q))`
# BST inorder = sorted; validate with (lo, hi) bounds
# serialize: preorder with "#" for None
\end{py}

% =====================================================================
\section{Templates: arrays and searching}

\subsection{Two pointers}
\begin{py}
l, r = 0, len(a) - 1                 # sorted a: pair sum
while l < r:
    s = a[l] + a[r]
    if s == t: ...; l += 1; r -= 1
    elif s < t: l += 1
    else: r -= 1
# fast/slow pointers: cycle detection, middle of list
# read/write pointers: in-place filtering
\end{py}

\subsection{Sliding window}
\begin{tpl}
def longest_window(a, ok_to_add, add, remove):
    # generic variable window: grow right, shrink left
    best, l = 0, 0
    for r, x in enumerate(a):
        while not ok_to_add(x):
            remove(a[l])
            l += 1
        add(x)
        best = max(best, r - l + 1)
    return best
\end{tpl}
Counting windows: if every window ending at \cd{r} with start \cd{>= l} is valid, add \cd{r - l + 1}.
``Exactly $k$'' $=$ ``at most $k$'' $-$ ``at most $k-1$''. Fixed size $k$: add \cd{a[r]}, drop \cd{a[r-k]}.

\subsection{Prefix sums}
\begin{tpl}
def count_subarrays_with_sum(nums, k):
    seen, pre, cnt = {0: 1}, 0, 0    # prefix -> frequency
    for x in nums:
        pre += x
        cnt += seen.get(pre - k, 0)  # look up before insert
        seen[pre] = seen.get(pre, 0) + 1
    return cnt

def prefix_2d(g):
    R, C = len(g), len(g[0])
    P = [[0] * (C + 1) for _ in range(R + 1)]
    for r in range(R):
        for c in range(C):
            P[r+1][c+1] = g[r][c] + P[r][c+1] + P[r+1][c] - P[r][c]
    return P
# sum of rows r1..r2, cols c1..c2 (inclusive):
# P[r2+1][c2+1] - P[r1][c2+1] - P[r2+1][c1] + P[r1][c1]
\end{tpl}
Variants: first index per prefix (longest subarray), key \cd{pre mod k} (divisible), parity/bitmask
prefix. \textbf{Difference array}: add $v$ on \cd{[l, r]} via \cd{D[l] += v; D[r+1] -= v}, then prefix-sum.

\subsection{Binary search on the answer}
\begin{tpl}
def max_true(lo, hi, ok):            # ok monotone True...False
    while lo < hi:
        mid = (lo + hi + 1) // 2     # round UP with lo = mid
        if ok(mid):
            lo = mid
        else:
            hi = mid - 1
    return lo

def min_true(lo, hi, ok):            # ok monotone False...True
    while lo < hi:
        mid = (lo + hi) // 2
        if ok(mid):
            hi = mid
        else:
            lo = mid + 1
    return lo
\end{tpl}
For floats, loop a fixed 100 times. Typical: capacity/speed/time minimisation, max side length, $k$-th
smallest.

\subsection{Monotonic stack}
\begin{tpl}
def next_greater(a):
    res, st = [-1] * len(a), []      # st holds indices
    for i, x in enumerate(a):
        while st and a[st[-1]] < x:
            res[st.pop()] = x
        st.append(i)
    return res
\end{tpl}
Also: largest rectangle in histogram, daily temperatures, sliding window max (monotonic deque).

\subsection{Intervals}
\begin{py}
iv.sort()                            # by start
out = []
for s, e in iv:
    if out and s <= out[-1][1]: out[-1][1] = max(out[-1][1], e)
    else: out.append([s, e])
# max non-overlapping: sort by END, keep if start >= last_end
# min rooms: sort starts and ends, or heap of end times
# sweep line: events (t, +1) / (t, -1), sort, running sum
\end{py}

\subsection{Greedy}
Make the locally best choice and prove it with an exchange argument, or stress-test it against brute
force. Common: farthest reach (jump games), earliest end (intervals), sort + take largest/smallest,
min/max heap of candidates.

% =====================================================================
\section{Templates: backtracking}
\begin{tpl}
def subsets(nums):
    res, path = [], []
    def bt(i):
        if i == len(nums):
            res.append(path[:])          # COPY
            return
        path.append(nums[i]); bt(i + 1); path.pop()   # take
        bt(i + 1)                                      # skip
    bt(0)
    return res

def combos(cands, target):          # reuse allowed
    cands, res, path = sorted(cands), [], []
    def bt(start, rem):
        if rem == 0:
            res.append(path[:]); return
        for i in range(start, len(cands)):
            if cands[i] > rem: break          # sorted -> prune
            path.append(cands[i])
            bt(i, rem - cands[i])             # i+1: no reuse
            path.pop()
    bt(0, target)
    return res

def perms(nums):
    res, path, used = [], [], [False] * len(nums)
    def bt():
        if len(path) == len(nums):
            res.append(path[:]); return
        for i, x in enumerate(nums):
            if not used[i]:
                used[i] = True; path.append(x)
                bt()
                path.pop(); used[i] = False
    bt()
    return res
\end{tpl}
\textbf{Duplicates in input:} sort, then skip \cd{i > start and a[i] == a[i-1]} (permutations: skip if
\cd{a[i] == a[i-1] and not used[i-1]}). \textbf{Grid paths:} mark the cell, recurse into neighbours, unmark.
\textbf{Pruning:} sort descending to fail fast; bound by the remaining sum; skip symmetric states (e.g.\ equal
bucket sums); check necessary conditions first (counts, divisibility).

% =====================================================================
\section{Templates: dynamic programming}
Steps: define the state in words, write the transition and the base case, pick an order so dependencies are
ready, then read off the answer. Top-down \cd{@cache} is quickest to write; bottom-up avoids recursion and is
faster.

\begin{tabular}{@{}p{2.4cm}p{5.4cm}@{}}
\toprule
Family & Core recurrence \\
\midrule
Kadane & \cd{cur = max(x, cur + x)} \\
House robber & \cd{dp[i] = max(dp[i-1], dp[i-2] + a[i])} \\
Grid paths & \cd{dp[r][c] = dp[r-1][c] + dp[r][c-1]} \\
0/1 knapsack & items outer, capacity \textbf{descending} \\
Unbounded & capacity \textbf{ascending} (or items inner) \\
Count combos & coins outer, amount inner \\
Count orders & amount outer, coins inner \\
String prefix & \cd{dp[i] = any(dp[j] and s[j:i] ok)} \\
LIS & tails + \cd{bisect_left}: $O(n \log n)$ \\
LCS / edit & \cd{dp[i][j]} from \cd{(i-1,j-1),(i-1,j),(i,j-1)} \\
Interval & \cd{dp[i][j]} by increasing length \\
Minimax game & \cd{f(i,j) = max(a[i]-f(i+1,j), a[j]-f(i,j-1))} \\
Bitmask & \cd{dp[mask]}, \cd{n <= 20} \\
State machine & \cd{dp[i][state]} (stock buy/sell/cooldown) \\
\bottomrule
\end{tabular}

\begin{tpl}
def knapsack_01(weights, values, cap):
    dp = [0] * (cap + 1)
    for w, v in zip(weights, values):
        for c in range(cap, w - 1, -1):     # descending: once
            dp[c] = max(dp[c], dp[c - w] + v)
    return dp[cap]

def min_coins(coins, amount):               # unbounded
    INF = amount + 1
    dp = [0] + [INF] * amount
    for a in range(1, amount + 1):
        for c in coins:
            if c <= a:
                dp[a] = min(dp[a], dp[a - c] + 1)
    return dp[amount] if dp[amount] < INF else -1

def lis_length(a):
    import bisect
    tails = []                  # tails[k] = min tail of an inc. seq of len k+1
    for x in a:
        i = bisect.bisect_left(tails, x)
        tails[i:i + 1] = [x]
    return len(tails)

def edit_distance(s, t):
    prev = list(range(len(t) + 1))
    for i in range(1, len(s) + 1):
        cur = [i] + [0] * len(t)
        for j in range(1, len(t) + 1):
            if s[i - 1] == t[j - 1]:
                cur[j] = prev[j - 1]
            else:
                cur[j] = 1 + min(prev[j - 1], prev[j], cur[j - 1])
        prev = cur
    return prev[-1]

def interval_game(a):           # first player's score, both optimal
    n, pre = len(a), [0]
    for x in a:
        pre.append(pre[-1] + x)
    dp = a[:]                   # dp[i] = best for player to move on a[i:i+L]
    for L in range(2, n + 1):
        dp = [pre[i+L] - pre[i] - min(dp[i], dp[i+1]) for i in range(n - L + 1)]
    return dp[0]
\end{tpl}
Space: if a row depends only on the previous row, keep two rows (or one, with the right loop direction).

% =====================================================================
\section{Simulation techniques}
\begin{itemize}
  \item \textbf{Synchronous updates:} read the old state only. Write to a new grid, or keep the next state
        in spare bits (\cd{cell |= 2}, then \cd{cell >>= 1}).
  \item \textbf{Deleting while scanning:} use an index-based \cd{while} loop; decide whether the index stays
        or advances after each deletion; keep parallel arrays in sync.
  \item \textbf{Rounds until stable:} loop until a round changes nothing (track a \cd{changed} flag).
  \item \textbf{Huge repetition:} detect a cycle of states (dict state $\to$ step), or track only
        offsets/extremes, or use math.
  \item \textbf{Robot/direction:} index into a clockwise list; positions as \cd{(r, c)}; validate each move.
  \item \textbf{Try every single edit:} precompute prefix states so each candidate only re-simulates the suffix.
  \item \textbf{Neighbourhood ops (blur, counts):} loop \cd{dr, dc in (-1, 0, 1)}, skip out of bounds, divide by the
        actual neighbour count; check the rounding rule (\cd{//}, round half up, banker's).
\end{itemize}

% =====================================================================
\section{Math and counting}
\begin{py}
MOD = 10**9 + 7                   # reduce after every op
# sieve of primes up to n
is_p = [True] * (n + 1); is_p[0] = is_p[1] = False
for i in range(2, math.isqrt(n) + 1):
    if is_p[i]:
        is_p[i*i::i] = [False] * len(range(i*i, n + 1, i))
# arithmetic series: n * (n + 1) // 2 ;  pairs: n * (n - 1) // 2
# count paths of length 2 through v: deg(v) * (deg(v) - 1) // 2
# count length-3 paths: fix the MIDDLE edge, multiply end choices
# Manhattan distance |dx| + |dy|; rotate with (x + y, x - y)
\end{py}

% =====================================================================
\section{Index mapping and gradients}
\begin{py}
order = sorted(range(n), key=lambda i: x[i])  # sorted pos -> orig idx
s = [x[i] for i in order]
grad_x = [0.0] * n
for k, i in enumerate(order):
    grad_x[i] = grad_s[k]       # a permutation's backward = scatter back
\end{py}
Chain rule over every path: a value used twice (e.g.\ an element that is also the median/normaliser)
receives the sum of both contributions. For $y = u / m$: $\partial y/\partial u = 1/m$,
$\partial y/\partial m = -u/m^2$. Check against a central difference
$\bigl(f(x + h) - f(x - h)\bigr)/2h$ with $h \approx 10^{-6}$.

% =====================================================================
\section{Testing your solution fast}
\begin{py}
import random
for _ in range(1000):                # stress test
    n = random.randint(1, 8)
    a = [random.randint(-5, 5) for _ in range(n)]
    assert fast(a) == brute(a), a     # prints the failing input
\end{py}
Hand-check: empty, one element, all equal, all negative, sorted/reversed, duplicates, max size (time), deepest
recursion, disconnected graph, no answer ($-1$/0/\cd{False}).

% =====================================================================
\section{Pre-submit checklist}
\begin{itemize}
  \item Output type and format exactly as asked (int vs.\ bool, string spacing, 0- vs.\ 1-indexed).
  \item Ties broken as specified; ``every'' vs.\ ``exists''; strict vs.\ non-strict comparisons.
  \item Hidden $O(n)$ in a loop: \cd{in list}, \cd{pop(0)}, slicing, \cd{sum}, string concatenation.
  \item Recursion depth over about 1000 frames?
  \item BFS marks visited on enqueue; backtracking appends copies.
  \item Inputs mutated that shouldn't be? Global state or caches reset between calls?
  \item Integer division and modulo semantics with negatives.
\end{itemize}

\end{multicols}
\end{document}
